\section{More detail on the three problems}
\label{app:failures}

This section gives the steps behind Section~\ref{sec:failures} that the main text
states briefly.

\subsection{Why making the candidates equally loud is not enough}
\label{app:conf}

Every step of the envelope pipeline is linear, so scaling a recording by $\alpha$
scales its envelope by $\alpha$: $\mathrm{env}(\alpha x)=\alpha\,\mathrm{env}(x)$.
Standardising each candidate then subtracts its mean and divides by its standard
deviation, which removes any change of the form $x\mapsto\alpha x+\beta$. Playing a
recording $15$\,dB louder or quieter therefore gives exactly the same
standardised envelope. A pure difference in loudness cannot survive, so it cannot
be what gives the answer away.

What does survive is every property that scaling and shifting leave unchanged.
Standardising fixes only two numbers per candidate, the mean and the variance.
Skew, kurtosis, sparsity, ratios between percentiles and the balance between
frequency bands all pass through untouched. If the attended and ignored
recordings differ in any of these, no rescaling can help --- and on this dataset
they do: the envelope's peak-to-average ratio alone separates attended and ignored
talkers with an effect size of $d=\R{sep/envelopecrestdB/d}$. A fix must therefore either make the candidates'
whole distributions identical (\cMat{}) or remove the difference between
attended and ignored recordings altogether (same-talker candidates). Reporting
chance does not help either, since the mark lifts accuracy well above it.

\subsection{How the standard score reduces to an audio classifier}
\label{app:degenerate}

The standard score has one convolution stack for the recording and one, shared
across candidates, for the audio. At each time step $t$ it normalises the brain
embedding $c_t$ and the candidate embedding $a_{k,t}$ to unit length, multiplies
them element by element, averages over time and applies a linear layer:
$\mathrm{score}_k = w^\top\big(\frac1T\sum_t \hat c_t\odot\hat a_{k,t}\big)+b$.
Now let the brain encoder output the same unit vector at every step,
$\hat c_t=\hat c$. The brain term no longer depends on $t$, so it moves out of the
average:
\begin{equation}
  \mathrm{score}_k \;=\; w^\top\!\Big(\hat c \odot \tfrac1T\textstyle\sum_t \hat a_{k,t}\Big) + b
  \;=\; \tilde w^\top \bar a_k + b,\qquad \tilde w=w\odot\hat c .
\end{equation}
This is a linear classifier on the time-averaged audio embedding $\bar a_k$ of
each candidate, and it can read any audio mark the audio encoder learns to
represent. The model reaches it simply by letting the brain encoder go constant,
which is easier than learning to follow the speech. The time-centred score of
Eq.~\ref{eq:coupling} closes this route: with a constant brain embedding,
$e-\bar e=0$, every correlation is zero, and every candidate gets the same score.

\subsection{Why a joint shuffle cannot see crowding out}
\label{app:freeload}

Shuffling all recorded streams together measures how much the model uses the
listener's signals as a whole. Suppose EEG alone earns a contribution of $0.40$.
Now add a gaze stream that is easier to read, and suppose the fused model earns
$0.55$ in total, but only $0.20$ of it from the EEG. The joint contribution went
\emph{up}, from $0.40$ to $0.55$, while the brain's share fell by half. A joint
shuffle reports the first change and hides the second. Only a split per stream,
such as the Shapley values of Eq.~\ref{eq:shapley}, shows both. The numbers of
Table~\ref{tab:hinge} follow exactly this pattern.

\subsection{Why we shuffle the recording, not the labels}
\label{app:perm}

A common check shuffles the labels and retrains or re-scores. That is the wrong
test here. Shuffling the labels breaks the link between the audio and the answer
as well as the link between the brain and the answer, so a decoder that reads
only the audio also drops to chance, and the shortcut we are testing for is
hidden. We instead shuffle the \emph{recording} across test windows and keep each
window's own candidates and label. A decoder that reads the audio keeps its
accuracy; only a decoder that uses the recording loses any. The accuracy that
remains, $v(\emptyset)$, is what \emph{this} model gets from the audio alone,
measured through its own audio pathway. The linear probe bounds what a simple
reader can take from the audio; $v(\emptyset)$ measures what this decoder
actually takes.
